\chapter{累次积分与变量代换}
\label{chap:fubini-change-variables-solutions}

\begin{solution}{fubini-uniform-section-proof}
设 \(f\) 在 \(Q=[a,b]\times[c,d]\) 上连续。给定 \(\varepsilon>0\)，由
一致连续性取 \(y\) 方向分割 \(P_y\)，使每个小区间内
\[
 |f(x,y)-f(x,y')|
 <\frac{\varepsilon}{2(b-a)(d-c)}
 \qquad(x\in[a,b])
\]
成立。于是任取样点 \(\eta_j\)，内层和满足一致估计
\[
 \left|\sum_j f(x,\eta_j)\Delta y_j-
       \int_c^d f(x,y)\,dy\right|
 <\frac{\varepsilon}{2(b-a)}
 \qquad(x\in[a,b]).
\]
把这个误差乘以外层小区间长度并求和，总误差小于
\(\varepsilon/2\)。再把 \(x\) 方向分割取得充分细，使 \(F(x)\) 的
Riemann 和与 \(\int_a^bF\) 相差小于 \(\varepsilon/2\)。同一个乘积
Riemann 和于是同时逼近多重积分与累次积分。紧致性用来把 \(f\) 的连续性
提升为对全部 \(x,y\) 同时有效的一致连续估计。
\end{solution}

\begin{solution}{triangle-iterated-integral}
令
\[
 D=\{(x,y):a\le x\le b,\ \alpha(x)\le y\le\beta(x)\},
\]
其中 \(\alpha,\beta\) 连续。两条连续图像在紧区间上一致连续。把
\([a,b]\) 分得足够细，每段图像都落在高度任意小的矩形条中，故两条图像
Jordan 零测；连同两端竖线可知 \(\partial D\) Jordan 零测。因此 \(D\)
Jordan 可测。

把 \(f\) 连续延拓到外包矩形 \(Q=[a,b]\times[c,d]\)。函数
\(f\mathbf1_D\) 只可能在 \(\partial D\) 上不连续，故可积。对每个 \(x\)，
\[
 \int_c^d f(x,y)\mathbf1_D(x,y)\,dy
 =\int_{\alpha(x)}^{\beta(x)}f(x,y)\,dy.
\]
Riemann 可积函数的累次积分定理给出
\[
 \int_Df
 =\int_a^b\int_{\alpha(x)}^{\beta(x)}f(x,y)\,dy\,dx.
\]
\end{solution}

\begin{solution}{riemann-section-nonintegrable}
在 \([-1,1]^2\) 上令
\[
 f(x,y)=
 \begin{cases}
  \mathbf1_{\Q}(y),&x=0,\\
  0,&x\ne0.
 \end{cases}
\]
若 \(x\ne0\)，函数在该点附近恒为零；若 \(x=0\)，每个邻域都同时含有
函数值 \(0\) 与 \(1\) 的点。因此不连续点集恰为竖线
\(\{0\}\times[-1,1]\)，这是 Jordan 零测集。振幅判据说明 \(f\)
Riemann 可积且积分为零。可是 \(x=0\) 的截面是 Dirichlet 函数
\(\mathbf1_{\Q}\)，在每个非退化区间上都不可 Riemann 积。
\end{solution}

\begin{solution}{fubini-riemann-oscillation}
设 \(P_A=(R_i)\)、\(P_B=(S_j)\)，并记
\[
 \omega_{ij}
 =\sup_{R_i\times S_j}f-\inf_{R_i\times S_j}f.
\]
若 \(x,x'\in R_i\)，则在每个 \(S_j\) 上
\[
 -\omega_{ij}\le f(x,y)-f(x',y)\le\omega_{ij}.
\]
两侧积分并对 \(j\) 求和，得到
\[
 |F(x)-F(x')|\le\sum_j\omega_{ij}|S_j|.
\]
取 \(x,x'\) 的上确界，乘以 \(|R_i|\) 后对 \(i\) 求和，即得
\[
 U(F,P_A)-L(F,P_A)
 \le U(f,P_A\times P_B)-L(f,P_A\times P_B).
\]
右端可任意小，所以只由 Darboux 判据就能推出 \(F\) 可积；证明不需要
\(F\) 连续。随后任取外层标记点，截面积分的上下和仍夹在同一乘积
Darboux 上下和之间，令其差趋零便得到积分等式。
\end{solution}

\begin{solution}{linear-change-volume}
平移不改变体积。对线性部分，坐标置换只重新排列累次积分的次序，反射只
反转某一截面的方向，二者均保持无向体积。对角伸缩
\[
 (x_1,\ldots,x_n)\longmapsto
 (\lambda_1x_1,\ldots,\lambda_nx_n)
\]
由一元变量代换把体积乘以 \(\prod_i|\lambda_i|\)。初等剪切
\[
 (x_1,\ldots,x_i,\ldots,x_n)
 \longmapsto
 (x_1,\ldots,x_i+c x_j,\ldots,x_n)
\]
在固定其余坐标后只平移第 \(i\) 个方向的截面，故由累次积分保持体积。
任意可逆矩阵是这些初等矩阵的乘积；各步体积倍率相乘，而各步行列式也
相乘，得到
\[
 |A(Q)|=|\det A|\,|Q|.
\]
同一论证适用于经过前面步骤得到的平行多面体，因为其相关截面仍是有限个
区间所成的 Jordan 可积集合。
\end{solution}

\begin{solution}{local-volume-distortion-details}
记 \(Q_0=[-r,r]^n\)，
\[
 g(h)=A^{-1}\bigl(\Phi(a+h)-\Phi(a)\bigr)-h.
\]
由正文估计，\(\norm{g(h)}_\infty\le\eta_rr\)。若
\(\norm y_\infty\le(1-\eta_r)r\)，则
\[
 \norm{T_y(h)}_\infty
 =\norm{y-g(h)}_\infty
 \le(1-\eta_r)r+\eta_rr=r,
\]
所以 \(T_y\) 把 \(Q_0\) 映入自身。其 Lipschitz 常数不超过
\(\eta_r<1\)，故有唯一不动点。

内外包给出
\[
 (1-\eta_r)^n|\det A|\,|Q_r|
 \le|\Phi(Q_r)|
 \le(1+\eta_r)^n|\det A|\,|Q_r|.
\]
因而相对误差满足
\[
 \frac{\bigl||\Phi(Q_r)|-|\det A|\,|Q_r|\bigr|}
      {|\det A|\,|Q_r|}
 \le
 \max\{(1+\eta_r)^n-1,\ 1-(1-\eta_r)^n\}.
\]
当 \(0\le\eta_r\le1/2\) 时，右端不超过
\(n2^{n-1}\eta_r\)，这可由中值定理作用于 \(t\mapsto(1+t)^n\) 得到。
\end{solution}

\begin{solution}{change-variable-orientation-absolute}
取 \(T(u,v)=(-u,v)\)。它把单位正方形映为面积仍为 \(1\) 的正方形，而
\(\det DT=-1\)。若变量代换公式使用 \(\det DT\) 而不用绝对值，常函数
\(1\) 的积分会变成 \(-1\)，与 Jordan 体积非负矛盾。行列式符号记录
定向，Riemann 体积只记录无向体积。
\end{solution}

\begin{solution}{jacobian-zero-no-local-inverse}
一维映射 \(f(x)=x^3\) 在全实轴严格单射，但 \(f'(0)=0\)。其逆映射
\(y\mapsto y^{1/3}\) 在零点没有有限导数，因而不是 \(C^1\) 局部微分
同胚。Jacobian 为零不排除集合意义上的单射，却使线性主部不可逆，不能
满足逆函数定理的假设。
\end{solution}

\begin{solution}{polar-coordinate-disk}
在 \(0<r<1\)、\(0<\theta<2\pi\) 上，极坐标 Jacobian 为 \(r\)，故
\[
 \int_{x^2+y^2\le1}(x^2+y^2)\,dx\,dy
 =\int_0^{2\pi}\int_0^1r^3\,dr\,d\theta
 =\frac{\pi}{2}.
\]
原点对应整条 \(r=0\)，切口两侧对应 \(\theta=0,2\pi\)。先在
\(\rho\le r\le1\)、\(\eta\le\theta\le2\pi-\eta\) 上应用变量代换，再令
\(\rho,\eta\downarrow0\)。删去部分及其像的体积趋于零，所以极限给出
圆盘上的公式。
\end{solution}

\begin{solution}{spherical-coordinate-jacobian}
取
\[
 \Phi(r,\varphi,\theta)
 =(r\sin\varphi\cos\theta,\,
   r\sin\varphi\sin\theta,\,
   r\cos\varphi).
\]
三个偏导列向量的混合积为
\[
 \det D\Phi=r^2\sin\varphi.
\]
在 \(r=0\) 和 \(\varphi=0,\pi\) 处退化；\(\theta=0\) 与 \(2\pi\)
表示同一半平面，故闭参数盒上非单射。应先在
\[
 r>0,\qquad0<\varphi<\pi,\qquad0<\theta<2\pi
\]
的相对紧子盒上使用定理，再由 Jordan 零测边界和积分估计补回极轴、
原点与切口。
\end{solution}

\begin{solution}{piecewise-diffeomorphism-multiplicity}
在参数域 \((0,1)\times(0,4\pi)\) 中，除切口外圆盘内每个非零点恰有
两个原像，分别相差 \(2\pi\)。把角区间分成两个周期并分别变量代换，对
连续 \(g\) 得
\[
 \int_0^{4\pi}\int_0^1
 g(r\cos\theta,r\sin\theta)r\,dr\,d\theta
 =2\int_{B(0,1)}g(x,y)\,dx\,dy.
\]
原点和切口 Jordan 零测，不影响等式。这正是重数函数在非例外点等于
\(2\) 的情形。
\end{solution}

\begin{solution}{riemann-change-variables-boundary-audit}
紧致性把 \(D\Phi\) 的连续性提升为全体小立方体上的一致误差，并给
\(D\Phi\)、\(D\Phi^{-1}\) 统一界；Jordan 边界使边界小块总体积趋于零；
微分同胚保证局部内包、全局不重复计数以及边界映成边界；非零 Jacobian
使局部线性映射可逆。绝对值则把有向体积改为普通体积。

删除全局单射时，两个角周期的极坐标使同一像点被计数两次。删除非退化性
时，映射 \((u,v)\mapsto(u,0)\) 把二维区域压到直线，不能再用同维体积
公式。若 \(U\) 的边界不是 Jordan 零测集，边界网格的总贡献未必趋于零，
本章的有限分割证明也会失去依据。
\end{solution}
